% Part: first-order-logic
% Chapter: syntax
% Section: intro-syntax

\documentclass[../../../include/open-logic-section]{subfiles}

\begin{document}

\olfileid{fol}{syn}{itx}

\olsection{Pambuka}

Supaya bisa ngrembakakake teori lan metateori logika sepisanan, kita
luwih dhisik kudu nemtokake sintaksis lan semantik ungkapan-ungkapane.
Ungkapan logika sepisanan iku term lan !!{formula}s. Term kawangun saka
!!{variable}s, !!{constant}s, lan !!{function}s. Sabanjure,
!!^{formula}s kawangun saka !!{predicate}s bebarengan karo term
(iki mbentuk !!{formula}s ``atomik'' kang paling prasaja), banjur saka
!!{formula}s atomik kita bisa mbentuk formula kang luwih rumit kanthi
nggunakake konektif logis lan kuantor. Ana akeh cara kanggo netepake
aturan pambentukan; ing kene kita mung menehi salah siji cara kang
bisa dipilih. Sistem liyane bakal milih simbol kang beda, milih
himpunan konektif kang beda minangka konektif primitif, lan nggunakake
tandha kurung kanthi cara kang beda (utawa malah ora nggunakake babar
pisan, kaya ing notasi kang diarani notasi Polandia). Cathetan koreksi
sumber OLPL-063: sumber Inggris nggunakake ``will chose'', yaiku wujud
kala kapungkur sawise ``will''; terjemahan iki mulihake wujud ``will
choose'' kang dikarepake. Nanging, kabeh pendekatan nduweni bab kang
padha: aturan pambentukan nemtokake himpunan term lan !!{formula}s
kanthi \emph{induktif}. Yen ditindakake kanthi trep, saben ungkapan
sejatine mung bisa kawangun kanthi siji cara miturut aturan pambentukan.
Definisi induktif kang ngasilake ungkapan kang \emph{mung bisa diwaca
kanthi siji cara} tegese kita bisa menehi makna marang ungkapan kasebut
kanthi nggunakake metode kang padha---definisi induktif.

\end{document}
